%%%%%%%%%%  scan idx 99  |  folio 86  %%%%%%%%%%
\SGAPrintPage{86}%
We shall say that \(\underline{C}_0\) is \uline{locally stable under kernels of epimorphisms} when the following condition is satisfied: for every \(X\in\Ob \underline{S}\) and every epimorphism \(u:E\longrightarrow F\) of \(\underline{C}_X\) with \(E\) and \(F\in\Ob\underline{C}_{0X}\), the kernel of \(u\) is locally in \(\underline{C}_0\) (i.e.\ the set of \(f:U\to X\) in \(\underline{S}_{/X}\) such that \(f^*\Ker(u)\in\Ob\underline{C}_{0U}\) is a refinement of \(X\)).

By the standard procedure of decomposing a long exact sequence into short exact sequences, we see that the preceding condition is equivalent to the following (apparently stronger) condition: for every \(X\in\Ob \underline{S}\) and every exact sequence of \(\underline{C}_X\)
\sgaStableTarget{sga6.e01.eq.unnamed}\begin{equation}\tag{*}
  0\longrightarrow E^0\longrightarrow E^1\longrightarrow\ldots\longrightarrow E^n\longrightarrow 0
\end{equation}
with \(E^i\in\Ob\underline{C}_{0X}\) for \(1\leq i\leq n\), \(E^0\) is locally in \(\underline{C}_0\).

We shall say that \(\underline{C}_0\) is \uline{locally quasi-stable under kernels of epimorphisms} when, for every exact sequence such as \((*)\), \(E^0\) is of \(\underline{C}_0\)-finite type.

On the other hand, we shall say that \(\underline{C}_0\) is \uline{locally liftable in} \(\underline{C}\) when, for every \(X\in\Ob \underline{S}\), every epimorphism \(u:E\longrightarrow F\) of \(\underline{C}_X\) with \(F\in\Ob\underline{C}_{0X}\) locally admits a section. For \(\underline{C}_0\) to be locally liftable in \(\underline{C}\), it is necessary and sufficient that, for every diagram of \(\underline{C}_X\) (\(X\in\Ob \underline{S}\))
\[
\begin{tikzcd}
& H \arrow[d]\\
E \arrow[r,"u"] & F
\end{tikzcd}
\]
where \(u\) is an epimorphism and \(H\in\Ob\underline{C}_0\), the set of \(U\) over \(X\) such that there exists a commutative diagram
\[
\begin{tikzcd}
& H_U \arrow[d] \arrow[dl]\\
E_U \arrow[r,"u_U"] & F_U
\end{tikzcd}
\]
is a refinement of \(X\).

%%%%%%%%%%  scan idx 100  |  folio 87  %%%%%%%%%%
\SGAPrintPage{87}%
We shall say that \(\underline{C}_0\) is \uline{locally quasi-liftable in} \(\underline{C}\) when the following weaker condition is satisfied: for every \(X\in\Ob \underline{S}\), and every epimorphism \(u:E\longrightarrow F\) of \(\underline{C}_X\) with \(F\in\Ob\underline{C}_{0X}\), the set of \(f:U\to X\) in \(\underline{S}_{/X}\) such that there exists \(v:F'\longrightarrow f^*E\) with \(F'\in\Ob\underline{C}_{0U}\) such that \(f^*(u)v\) is an epimorphism is a refinement of \(X\).

This condition is equivalent to the following: for every diagram of \(\underline{C}_X\)
\[
\begin{tikzcd}
& H \arrow[d]\\
E \arrow[r,"u"] & F ,
\end{tikzcd}
\]
where \(u\) is an epimorphism, and \(H\in\Ob\underline{C}_0\), the set of \(U\) over \(X\) such that there exists a commutative diagram
\[
\begin{tikzcd}
G \arrow[r,"v"] \arrow[d] & H_U \arrow[d]\\
E_U \arrow[r,"u_U"] & F_U ,
\end{tikzcd}
\]
with \(G\in\Ob\underline{C}_0\) and \(v\) an epimorphism is a refinement of \(X\).

When \(\underline{S}\) is a chaotic site, for example the punctual site (one object, one morphism), we shall omit the adverb ``locally'' in the preceding definitions.

\sgaStableTarget{sga6.e01.sec.1.3}\subsection*{1.3. Examples}

a) Let \(\underline{S}\) be a punctual site, \(A\) a ring, and \(\underline{C}\) the category of left \(A\)-modules. The full subcategory \(\underline{C}_0\) of \(\underline{C}\) formed by the projective (left) \(A\)-modules of finite type is stable under kernels of epimorphisms and liftable in \(\underline{C}\); an \(A\)-module is of \(\underline{C}_0\)-finite type if and only if it is of finite type in the
%%%%%%%%%%  scan idx 101  |  folio 88  %%%%%%%%%%
\SGAPrintPage{88}%
ordinary sense. The full subcategory \(\underline{C}_1\) of \(\underline{C}\) formed by the \(A\)-modules of finite type is quasi-liftable in \(\underline{C}\) (but not liftable in general); when \(A\) is noetherian, \(\underline{C}_1\) is stable under kernels of epimorphisms. The full subcategory \(\underline{C}'_0\) of \(\underline{C}\) formed by the free \(A\)-modules of finite type is liftable in \(\underline{C}\) and quasi-stable (not stable in general) under kernels of epimorphisms.

b) Let \((S,\underline{O}_S)\) be a space ringed by local rings, \(\underline{S}\) the site of open subsets of \(S\), and \(\underline{C}\) the fibered \(\underline{S}\)-category (even split (SGA 1 VI 9)) defined by the functor \(U\longmapsto\) (category of left \(\underline{O}_U\)-Modules); \(\underline{C}\) is a flat abelian \(\underline{S}\)-category (1.1). Let \(\underline{C}_0\) be the strictly full additive sub-\(\underline{S}\)-category of \(\underline{C}\) defined by \(U\longmapsto\) (category of free \(\underline{O}_U\)-Modules of finite type). An \(\underline{O}_S\)-Module \(M\) is locally of \(\underline{C}_0\)-finite type if and only if \(M\) is of finite type in the ordinary sense (EGA \(0_{\mathrm I}\) 5.2). It follows from (EGA \(0_{\mathrm I}\) 5.2.7 and 5.4.9) that \(\underline{C}_0\) is locally stable under kernels of epimorphisms and locally liftable in \(\underline{C}\).

c) Let \(G\) be a group, \(\underline{S}\) the topos of \(G\)-sets (SGA 4 V 3.11), \(k\) a field (regarded as a constant Ring of \(\underline{S}\)), and \(\underline{C}\) the split \(\underline{S}\)-category defined by the functor \(U\longmapsto\) (category of left \(k_U\)-Modules); \(\underline{C}\) is a flat abelian \(\underline{S}\)-category. Let \(\underline{C}_0\) be the strictly full additive sub-\(\underline{S}\)-category of \(\underline{C}\) defined by \(U\longmapsto\) (category of free \(k_U\)-Modules of finite type).
Since localization in \(\underline{S}\) is equivalent to forgetting the operations of \(G\), it is clear that a \(G\)-\(k\)-module \(M\) is locally of \(\underline{C}_0\)-finite type if and only if \(M\) is finite-dimensional as a vector space over \(k\), and that \(\underline{C}_0\) is locally stable under kernels of epimorphisms and locally liftable in \(\underline{C}\).

The reader is left to examine the ``Galois'' variant of c) in which \(G\) is a profinite group and \(\underline{S}\) the topos of sets on which \(G\) operates continuously.

%%%%%%%%%%  scan idx 102  |  folio 89  %%%%%%%%%%
\SGAPrintPage{89}%
In fact, examples b) and c) are special cases of the following example:

d) Let \((\underline{S},\underline{O}_{\underline{S}})\) be a ringed topos (cf.\ SGA 4 IV 2.10)\SGAAuthorFootnote{*}{Here \(\underline{O}_{\underline{S}}\) is a ring of \(\underline{S}\), with identity element, but not necessarily commutative.}, and \(\underline{C}\) the split \(\underline{S}\)-category defined by \(U\longmapsto\) (category of left \(\underline{O}_U\)-Modules); \(\underline{C}\) is a flat abelian \(\underline{S}\)-category. Let \(\underline{C}_0\) be the strictly full additive sub-\(\underline{S}\)-category of \(\underline{C}\) defined by \(U\longmapsto\) (category of free \(\underline{O}_U\)-Modules of finite type). An \(\underline{O}_{\underline{S}}\)-Module locally of \(\underline{C}_0\)-finite type will simply be said to be of \uline{finite type}. The category \(\underline{C}_0\) is locally quasi-stable under kernels of epimorphisms and locally liftable in \(\underline{C}\), as follows from

\sgaStableTarget{sga6.e01.lem.1-3-1}\textbf{Lemma 1.3.1.} \uline{Let} \(P\) \uline{be an} \(\underline{O}_{\underline{S}}\)\uline{-Module. Equivalent conditions}:
\begin{enumerate}[label=(\roman*)]
\item \(P\) \uline{is of finite type and the functor} \(\underline{\Hom}(P,.)\) \uline{is exact};
\item \(P\) \uline{is of finite type and every epimorphism} \(M\to P\) \uline{locally admits a section};
\item \(P\) \uline{is locally a direct factor of a free Module of finite type}.
\end{enumerate}
The proof is immediate and is left as an exercise to the reader.

\sgaStableTarget{sga6.e01.sec.1.4}\subsection*{1.4. An ``extension'' lemma (cf. EGA \(0_{\mathrm{III}}\) 11.9.1)}

\sgaStableTarget{sga6.e01.prop.1-4-1}\textbf{Proposition 1.4.1.} \uline{Let} \(\underline{S}\) \uline{be a site,} \(\underline{C}\) \uline{a flat abelian} \(\underline{S}\)\uline{-category (1.1),} \(\underline{C}_0\) \uline{a strictly full additive sub-}\(\underline{S}\)\uline{-category (1.1) of} \(\underline{C}\)\uline{; assume that} \(\underline{C}_0\) \uline{is locally quasi-liftable in} \(\underline{C}\) \uline{(1.2). Let} \(S\in\Ob \underline{S}\)\uline{,} \(u:E\longrightarrow F\) \uline{be a morphism (of degree} \(0\)\uline{) of complexes of} \(\underline{C}_S\) \uline{(differentials of degree} \(+1\)\uline{),} \(G\) \uline{the cone of} \(u\) \uline{(cf.\ [V] chap.\ I \S~1 2.1),} \(n\in Z\)\uline{; assume that} \(\Ho^n(G)\) \uline{is}
%%%%%%%%%%  scan idx 103  |  folio 90  %%%%%%%%%%
\SGAPrintPage{90}%
\uline{of} \(\underline{C}_0\)-%
\uline{finite type (1.2). Then the set of} \(X\in\Ob \underline{S}_{/S}\) \uline{satisfying property} \((P)\) \uline{below is a refinement of} \(S\).

\((P)\): there exist an object \(M\) of \(\underline{C}_{0X}\) and arrows of \(\underline{C}_X\)
\[
  r:M\longrightarrow Z^{n+1}(E_X)=
  \Ker\!\left(E_X^{n+1}\xrightarrow{d^{n+1}}E_X^{n+2}\right)
\]
\[
  \text{and}\quad s:M\longrightarrow F_X^n
\]
such that the following diagram of \(\underline{C}_X\)
\[
\begin{tikzcd}[column sep=tiny]
\ldots \arrow[r] & E_X^{n-1}
 \arrow[r,"{\binom{d_{E_X}^{n-1}}{0}}"]
 \arrow[d,"u_X^{n-1}"] &
E_X^n\oplus M
 \arrow[r,"{(d_{E_X}^n,r)}"]
 \arrow[d,"{(u_X^n,s)}"] &
E_X^{n+1} \arrow[r] \arrow[d,"u_X^{n+1}"] &
\ldots\\
\ldots \arrow[r] & F_X^{n-1} \arrow[r] &
F_X^n \arrow[r] &
F_X^{n+1} \arrow[r] &
\ldots
\end{tikzcd}
\]
is a morphism of complexes \(u':E'\to F_X\) (i.e.\ \(d_{F_X}^n s=u_X^{n+1}r\)), and that
\begin{enumerate}[label=(\roman*)]
\item \(\Ho^{n+1}(u')\) is a monomorphism, with \(\Im\Ho^{n+1}(u')=\Im\Ho^{n+1}(u_X)\),
\item \(\Ho^n(u')\) is an epimorphism.
\end{enumerate}

In pictorial terms, we may say that, by modifying \(E\) and \(u\) locally, \uline{in degree \(n\) only}, by adding to \(E^n\) an object of \(\underline{C}_0\), we can locally make \(\Ho^{n+1}(u)\) injective, without changing its image, and \(\Ho^n(u)\) surjective.

\uline{Proof}. The procedure is as follows: first we show that we can choose \(r\) and \(s\) locally so as to make \(\Ho^{n+1}(u)\) injective; we verify 
